\documentclass[../cours_tle.tex]{subfiles}
\begin{document}
\setcounter{chapter}{32}
\chapter{Puissances de matrices et suites couplées}
\label{ch:33}
\begin{objectif}
Mathématiques expertes. Calculer une évolution répétée et identifier ses modes. Prérequis : chapitres 10 et 32.
\end{objectif}
\section{Définir une puissance}
Pour une matrice carrée, $A^0=I$ et $A^{n+1}=AA^n$. L'associativité donne $A^{n+m}=A^nA^m$.
Si $A=PDP^{-1}$ avec $D$ diagonale, alors $A^n=PD^nP^{-1}$ : au produit, les facteurs $P^{-1}P$ s'annulent ; une récurrence formalise ce calcul.
Trouver systématiquement une telle décomposition relève de la diagonalisation en licence ; ici on la vérifie lorsqu'elle est donnée ou accessible.
\section{Un calcul complet d'ordre 2}
Pour $A=\begin{pmatrix}2&1\\1&2\end{pmatrix}$, les vecteurs $(1,1)$ et $(1,-1)$ sont multipliés respectivement par $3$ et $1$.
Avec $P=\begin{pmatrix}1&1\\1&-1\end{pmatrix}$ et $P^{-1}=P/2$, on obtient
\[
A^n=\frac12\begin{pmatrix}3^n+1&3^n-1\\3^n-1&3^n+1\end{pmatrix}.
\]
On peut vérifier cette formule au rang $0$ puis par multiplication, sans invoquer un théorème de diagonalisation.
Un vecteur non nul $v$ tel que $Av=\lambda v$ est appelé vecteur propre ; la direction ne change pas, seul le facteur change.
\section{Matrices nilpotentes et puissances d'ordre 3}
Si $A=I+N$ avec $N^3=0$, alors
$A^n=I+nN+\binom n2N^2$.
Le binôme s'applique parce que $I$ et $N$ commutent ; les termes de degré au moins $3$ sont nuls.
Pour $N=\begin{pmatrix}0&1&0\\0&0&1\\0&0&0\end{pmatrix}$, cela donne
$A^n=\begin{pmatrix}1&n&n(n-1)/2\\0&1&n\\0&0&1\end{pmatrix}$.
Ce calcul montre qu'une puissance n'est pas obtenue en élevant séparément chaque coefficient à la puissance $n$.
\section{Suites affines de vecteurs}
Pour $U_{n+1}=AU_n+C$, une récurrence donne
$U_n=A^nU_0+\sum_{k=0}^{n-1}A^kC$.
Si $I-A$ est inversible, le point fixe est $L=(I-A)^{-1}C$ et
$U_n-L=A^n(U_0-L)$.
La convergence dépend donc des puissances de $A$, pas seulement de ses coefficients pris isolément.
Pour une matrice diagonale dont les coefficients ont des modules strictement inférieurs à $1$, chaque coordonnée converge vers le point fixe. Une décomposition $A=PDP^{-1}$ permet la même conclusion lorsque les coefficients diagonaux de $D$ satisfont cette condition.
\section{Récurrences d'ordre 2 et Fibonacci}
La relation $u_{n+2}=au_{n+1}+bu_n$ devient
$\binom{u_{n+2}}{u_{n+1}}=\begin{pmatrix}a&b\\1&0\end{pmatrix}\binom{u_{n+1}}{u_n}$.
Si $r^2-ar-b=0$ possède deux racines distinctes $r_1,r_2$, les suites $r_1^n,r_2^n$ sont solutions. Les conditions initiales déterminent une unique combinaison $u_n=\alpha r_1^n+\beta r_2^n$ ; l'unicité vient de la récurrence et des deux premières valeurs.
Pour Fibonacci $F_0=0,F_1=1$, les racines sont $\varphi=(1+\sqrt5)/2$ et $\psi=(1-\sqrt5)/2$, d'où
$F_n=(\varphi^n-\psi^n)/\sqrt5$.
On vérifie les deux valeurs initiales et la relation ; comme $|\psi|<1<\varphi$, le terme dominant est $\varphi^n/\sqrt5$.
Dans le cas d'une racine double non nulle $r$, les suites $r^n$ et $nr^n$ conviennent ; si la racine double est $0$, la relation impose $u_{n+2}=0$ et doit être traitée directement.
\begin{finchapitre}
\section{Utilité et suite du parcours}
Les modes propres expliquent croissance, oscillation et stabilité d'un système. Ils seront centraux en algèbre linéaire, équations différentielles et calcul scientifique.
\begin{retenir}
Pour une formule de puissance ou de suite, vérifier le rang initial et la relation. Une base de vecteurs propres n'existe pas toujours.
\end{retenir}
\end{finchapitre}
\end{document}
